\section{Limitations}

This manuscript reports internal validation, not expert-reviewed legal validity.
The live validation run is small and should be interpreted as an end-to-end
model-generation path test. The natural-response Dasha audit confirms that the
clustered inputs are model answers to the same Frank question rather than
responses generated from preassigned labels, but the current live batch is too
small to prove broad clustering robustness. The 500-response scale result is a
controlled regression test with known reasoning signatures, not a claim that
arbitrary live model outputs will always produce perfectly separable clusters.
The evidence is therefore strong for pipeline mechanics and internal
reproducibility, but still bounded for external legal validity and cross-doctrine
generalization.

The Statute-of-Frauds evidence should be read as a strong litmus test, not as an
exhaustive doctrinal proof. The live run validates one source-to-ranking case,
and the regression suite covers eight major Statute-of-Frauds reasoning
archetypes. A claim that the system handles every Statute-of-Frauds fact pattern
would require a larger held-out case suite across jurisdictions, transaction
types, writing forms, and exception doctrines.

The current live model roster is intentionally small. The manuscript validates
the path with \texttt{gpt-5.2}, \texttt{claude-sonnet-4-20250514}, and
\texttt{meta/meta-llama-3-70b-instruct}; broader model-family claims require
larger rosters, repeated samples, held-out cases, and natural-response batches
where multiple models independently converge on or diverge from the same legal
reasoning paths.

The instruction corpus is deepest for Statute of Frauds. The live agents are now
given a doctrine-general protocol and must infer doctrine from the source, but
other doctrine families will need richer pack-specific examples before external
claims should be made about broad legal coverage.

LLM-as-judge methods are known to have bias risks, including position,
verbosity, and self-preference biases
\cite{zheng2023judging,wang2024positionbias,panickssery2024selfpreference}. The
pipeline mitigates some of these risks by using row-level rubrics, source
support, centroid-level scoring, and model-output provenance, but future
validation should include judge ensembles, repeated scoring, rubric ablation,
model-separated judges, and comparison to expert labels.
